
% !TeX root = ../main.tex
\begin{chineseabstract}
% \footnotetext{*本研究得到某某基金（编号：）的资助}

基于 Wi-Fi 信道状态信息（Channel State Information，CSI）的室内动作识别能够利用人体运动引起的信道变化实现非视觉感知，在智能家居、健康监护和环境感知等场景中具有应用价值。与固定类别条件下的离线识别不同，实际系统中的动作类别往往会随应用需求逐步增加，同时旧类数据难以长期完整保留，新类样本数量也可能不足。传统依赖全量样本集中训练的方法难以适应此类持续更新需求。针对上述问题，本文围绕 Wi-Fi CSI 室内动作识别中的类别持续扩展问题，分别研究新增类样本相对充足的类增量学习和新增类样本极少的少样本类增量学习两类场景。

针对类增量学习场景，本文提出一种回放—扩展—巩固框架。该方法通过记忆库重放保留旧类代表样本，利用冻结骨干与分类器扩展实现新类别接入，并结合教师—学生结构与双边知识蒸馏完成新旧类知识巩固，以缓解灾难性遗忘和新旧类分布失衡。针对少样本类增量学习场景，本文提出一种基于时序增强与图注意力的学习方法。该方法利用多尺度时序编码和时序注意力融合增强 CSI 动作时序表征，在增量阶段通过增强图注意力网络建模类别原型关系，并结合双分类器架构缓解新旧类偏置和类别级过拟合。

在 XRF 与 WiAR 两个公开 Wi-Fi 动作识别数据集上的实验结果表明，所提类增量学习方法在多阶段类别扩展过程中具有较好的增量稳定性。其中，在 XRF 数据集上末次增量阶段准确率达到 73.57\%；在 WiAR-16 类增量设置下，平均准确率达到 77.33\%，末次阶段准确率达到 64.58\%。所提少样本类增量学习方法在 XRF 的 5-way 1-shot 设定下末次增量会话准确率达到 74.24\%，平均准确率达到 78.48\%；在 WiAR 数据集上平均准确率为 75.30\%，末次会话准确率为 61.81\%，整体优于 CEC、FACT、C-FSCIL、DyCR 等对比方法。综上，本文从类增量学习和少样本类增量学习两个层面提出面向 Wi-Fi CSI 动作识别的持续学习方法，为动态场景下 Wi-Fi 感知系统的持续更新提供了可行思路。

\chinesekeywordstype{Wi-Fi CSI；室内动作识别；持续学习；类增量学习；少样本类增量学习}{应用研究}

\end{chineseabstract}

